%=====================================================================
%  VI_QLearning_explained.tex
%  A beginner-friendly explanation of Value Iteration (VI) and
%  Q-Learning IN THE CONTEXT OF THIS PROJECT:
%  "AI-Powered Autonomous Smart Agriculture using Swarm Intelligence
%   and Reinforcement Learning".
%
%  Compile with:  pdflatex VI_QLearning_explained.tex
%=====================================================================
\documentclass[11pt]{article}

\usepackage[a4paper,margin=1in]{geometry}
\usepackage{amsmath}
\usepackage{amssymb}
\usepackage{booktabs}
\usepackage{xcolor}
\usepackage{tcolorbox}
\usepackage{enumitem}
\usepackage{hyperref}
\hypersetup{colorlinks=true, linkcolor=blue!60!black, urlcolor=blue!60!black}

% Coloured callout boxes (same style as the PSO/GA guide).
\newtcolorbox{intuition}[1][]{colback=blue!4!white, colframe=blue!45!black,
  fonttitle=\bfseries, title=Intuition, #1}
\newtcolorbox{codebox}[1][]{colback=gray!6!white, colframe=gray!55!black,
  fonttitle=\bfseries, title=In our code, #1}
\newtcolorbox{whybox}[1][]{colback=green!4!white, colframe=green!40!black,
  fonttitle=\bfseries, title=Why we did this, #1}

\title{\textbf{Reinforcement Learning for Autonomous Navigation}\\[4pt]
\large A Plain-Language Guide to Value Iteration and Q-Learning\\ in the
Smart-Agriculture Project}
\author{AI Laboratory Project}
\date{}

\begin{document}
\maketitle

\begin{abstract}
\noindent
This document explains, in simple language, the two \emph{reinforcement
learning} algorithms used in Phase~2 of the project: \textbf{Value Iteration
(VI)} and \textbf{Q-Learning}. Both algorithms solve the \emph{same} question ---
\textbf{how should one robot move through the greenhouse to reach all of its
assigned flowers, avoid obstacles, and return safely to a charger?} --- but in
two very different ways. VI assumes the robot \emph{knows} the whole map
(model-based); Q-Learning assumes the robot knows \emph{nothing} and must learn
by trial and error (model-free). We describe the problem, the idea behind each
algorithm, every formula we use, and \emph{where and why} each design choice
appears in our code. No prior knowledge of reinforcement learning is assumed.
\end{abstract}

\tableofcontents

%=====================================================================
\section{The Problem We Are Solving (Phase 2: Navigation)}
%=====================================================================

In Phase~1, PSO (and the Genetic Algorithm) already decided \textbf{which flowers each robot
owns}. Phase~2 asks a completely different question, for \emph{one} robot at a
time:

\begin{center}
\emph{``How should this robot actually move, step by step, to reach all its
assigned flowers, avoid crashing into obstacles, use as little battery as
possible, and return safely to a charging station?''}
\end{center}

To make this concrete and solvable, we lay a \textbf{$20\times20$ grid} over the
greenhouse. The robot always sits in one grid cell and can take one of four
moves: \textbf{Up, Down, Left, Right}. Some cells contain obstacles (walls),
some contain the robot's target flowers, and some are charging stations.

\begin{intuition}
Think of it like a board game (or a phone maze game). The robot is a game piece
on a $20\times20$ board. Each turn it steps one square. We want to teach it the
\emph{smartest} sequence of steps to collect all its flowers and reach a
charger. ``Value Iteration'' and ``Q-Learning'' are two ways to figure out that
smart strategy.
\end{intuition}

\subsection{Why this is a Markov Decision Process (MDP)}

The navigation task has all the ingredients of a \textbf{Markov Decision
Process} --- the standard mathematical model for step-by-step decision making:

\begin{center}
\renewcommand{\arraystretch}{1.3}
\begin{tabular}{@{}p{2.6cm} p{11cm}@{}}
\toprule
\textbf{MDP part} & \textbf{In our project} \\
\midrule
\textbf{States} $s$ & Which grid cell the robot is in (row, column). \\
\textbf{Actions} $a$ & Up / Down / Left / Right (4 choices). \\
\textbf{Transitions} $T(s,a,s')$ & Where the robot ends up after an action ---
\emph{stochastic} (see below). \\
\textbf{Reward} $R$ & Points for good outcomes, penalties for bad ones. \\
\textbf{Discount} $\gamma$ & How much we value future reward vs immediate reward. \\
\bottomrule
\end{tabular}
\end{center}

\textbf{``Markov''} simply means: what happens next depends only on \emph{where
the robot is now} and \emph{what it does now} --- not on the whole history of how
it got there. This ``memoryless'' property is what makes the problem cleanly
solvable.

\subsection{Movement is stochastic (the robot can slip)}

A real robot's motors are not perfect. So when the robot \emph{intends} to go
(say) Up, we model the outcome as:

\begin{equation}
\begin{aligned}
&\text{go the intended way} && \text{with probability } 0.8,\\
&\text{slip to one perpendicular side} && \text{with probability } 0.1,\\
&\text{slip to the other perpendicular side} && \text{with probability } 0.1.
\end{aligned}
\end{equation}

\begin{intuition}
``I \emph{try} to step Up, but 20\% of the time a slippery floor pushes me
sideways.'' Because of this uncertainty, a good strategy must be \emph{robust}:
it should still work even when moves occasionally go wrong. This is why we cannot
just hard-code a single fixed path --- we need a \emph{policy} that says what to
do from \emph{every} cell, in case the robot slips into an unexpected one.
\end{intuition}

\begin{codebox}
Constants \texttt{MOVE\_PROB=0.8}, \texttt{SLIP\_PROB=0.1} (twice). The
\texttt{SLIP} dictionary maps each action to its two perpendicular slip
directions. Both VI and Q-Learning use these \emph{same} numbers.
\end{codebox}

\subsection{The reward design (how we tell the robot what ``good'' means)}

We shape the robot's behaviour purely through rewards and penalties:

\begin{center}
\renewcommand{\arraystretch}{1.3}
\begin{tabular}{@{}l r p{6.5cm}@{}}
\toprule
\textbf{Event} & \textbf{Reward} & \textbf{Why this value} \\
\midrule
Reach a flower (pollinate) & $+10$ & Big reward --- this is the main goal. \\
Reach charger (all flowers done) & $+10$ & Big reward --- safe return matters. \\
Each ordinary step & $-0.1$ & Small cost --- discourages wandering, saves battery/time. \\
Hit an obstacle/wall & $-5$ & Sharp penalty --- teaches collision avoidance. \\
\bottomrule
\end{tabular}
\end{center}

\begin{whybox}
The step cost of $-0.1$ is deliberately \emph{small but negative}. If it were
$0$, the robot would have no reason to hurry; if it were too large, the robot
might be too scared to explore. A gentle negative nudge makes the shortest safe
route the most rewarding one. The collision penalty ($-5$) is much larger than a
step cost so the robot strongly learns to avoid walls.
\end{whybox}

\begin{tcolorbox}[colback=yellow!8!white,colframe=orange!60!black,title=Note on convention: here, higher = better]
In Phase~1 (PSO/GA) we \emph{minimized} a cost (lower was better). In Phase~2
we \emph{maximize} a reward (higher is better). Same underlying idea, opposite
sign --- reinforcement learning is traditionally written in ``collect as much
reward as possible'' form.
\end{tcolorbox}

\subsection{The discount factor $\gamma$}
Reward that comes \emph{later} is worth a little less than reward \emph{now}.
The discount $\gamma\in[0,1)$ controls this: a reward $k$ steps in the future is
worth $\gamma^k$ times its face value. A $\gamma$ close to $1$ makes the robot
\emph{far-sighted} (willing to take many steps now for a big future reward),
which is exactly what we want for reaching a distant flower.

%=====================================================================
\section{What Are We Trying to Compute? Values and Policies}
%=====================================================================

Two central ideas appear in both algorithms:

\begin{itemize}[leftmargin=1.4em]
  \item \textbf{Value} $U(s)$ (or $V(s)$): ``Starting from cell $s$ and acting
        optimally forever after, how much total (discounted) reward can I
        expect?'' A cell right next to a flower has a high value; a cell in a
        dead-end far from everything has a low value.
  \item \textbf{Policy} $\pi(s)$: ``From cell $s$, which action should I take?''
        This is the actual \emph{strategy} --- a map giving the best move for
        every cell. The policy is what the robot ultimately follows.
\end{itemize}

\begin{intuition}
The \textbf{value} is like a ``heat map'' of how good each cell is. The
\textbf{policy} is the set of arrows telling you which way to walk from each
cell (always uphill on the value heat map). Both VI and Q-Learning are just two
recipes for finding these.
\end{intuition}

%=====================================================================
\section{Algorithm 1 --- Value Iteration (Model-Based)}
%=====================================================================

\subsection{The core idea}
Value Iteration assumes the robot \textbf{already knows the whole map}: where
every obstacle is, where the flowers are, and the exact slip probabilities. This
is the \textbf{model-based} setting --- we \emph{have} the transition model
$T(s,a,s')$. Given full knowledge, we can \emph{compute} the optimal values
directly, without the robot ever moving in the real world.

\begin{intuition}
``I have a perfect map and I know my motors slip 20\% of the time, so let me sit
down and \emph{calculate} the best plan before taking a single step.'' That is
Value Iteration: planning with complete information.
\end{intuition}

\subsection{The Bellman update (the heart of VI)}

We repeatedly refine our value estimate for every cell using the
\textbf{Bellman optimality update}:

\begin{equation}
U(s) \;\leftarrow\; \max_{a}\; \sum_{s'} T(s,a,s')\,\bigl[\,R(s,a,s') \;+\; \gamma\, U(s')\,\bigr].
\label{eq:bellman}
\end{equation}

Reading this in plain words, for one cell $s$:
\begin{enumerate}[nosep]
  \item For each of the 4 actions $a$, look at the cells you \emph{might} land in
        ($s'$), weighted by how likely each is ($T(s,a,s')=0.8/0.1/0.1$).
  \item For each possible landing, add the \textbf{immediate reward for that
        move}, $R(s,a,s')$: a normal step earns the step cost ($-0.1$); a move
        that gets blocked by a wall earns the collision penalty ($-5$). Then add
        the (discounted) value of where you land, $\gamma\,U(s')$.
  \item Averaging those over the $0.8/0.1/0.1$ outcomes gives ``the expected
        value of doing action $a$.''
  \item $\max_a$ --- keep the \emph{best} action's value (act optimally).
\end{enumerate}

Putting the reward \emph{inside} the sum (rather than a single $R(s)$ out front)
matters for one reason: it lets Value Iteration feel the \textbf{same collision
penalty} that Q-Learning feels whenever a move bumps a wall. Both algorithms
therefore optimise the \emph{identical} reward constants --- only the solving
method differs.

We apply equation~\eqref{eq:bellman} to \emph{every} cell, over and over. Each
sweep, the ``goodness'' of the flowers spreads outward one ring of cells at a
time, like ripples in a pond, until every cell knows how to get to a flower.

\begin{intuition}
Value ``flows'' outward from the flowers. After 1 sweep, only cells touching a
flower know they are good. After 2 sweeps, their neighbours know too. Keep
sweeping and eventually every cell has a correct value pointing the way to the
nearest reward. This is why the figure is called \emph{value propagation}.
\end{intuition}

\subsection{When do we stop? (convergence)}
Each sweep changes the values a little less than the last. We measure the biggest
change across all cells, $\delta$, and stop when it falls below a tiny threshold:
\[
\delta = \max_s \big| U_{\text{new}}(s) - U_{\text{old}}(s) \big| < \theta,
\qquad \theta = 10^{-4}.
\]
At that point the values have essentially stopped moving --- they have
\textbf{converged} to the true optimal values. In our run this happens in about
\textbf{33 iterations}.

\begin{codebox}
Function \texttt{value\_iteration(gw, goals, gamma=0.9, theta=1e-4)} in
\texttt{ValueIteration.py}. It is \emph{vectorized} with NumPy: instead of
looping cell by cell, the helper \texttt{neighbour\_value} shifts the whole value
grid in one direction at once, so all cells are updated together (fast and
exact). Goal cells are held fixed at $+10$ (they are \emph{terminal}); obstacle
cells are kept inert.
\end{codebox}

\subsection{From values to a policy}
Once the values have converged, the best move from any cell is simply the action
whose expected next-value is highest:
\begin{equation}
\pi(s) = \arg\max_{a}\; \sum_{s'} T(s,a,s')\,U(s').
\end{equation}
This turns the value ``heat map'' into a grid of arrows the robot can follow.

\begin{codebox}
Function \texttt{extract\_policy(gw, U, goals)} produces the arrow grid you see
in the policy figure. \texttt{run\_mission(...)} then plays the policy out:
because the robot must visit \emph{several} flowers, it re-plans after each
flower is reached (treat the next unvisited flower as the new goal), which is a
clean way to chain single-goal Value Iteration into a full multi-flower tour.
\end{codebox}

%=====================================================================
\section{Algorithm 2 --- Q-Learning (Model-Free)}
%=====================================================================

\subsection{The core idea}
Q-Learning throws away the map. Now the robot knows \textbf{nothing} in advance
--- no obstacle locations, no slip probabilities. It can only \emph{act}, then
\emph{observe} what happened (where it ended up, what reward it got). By
repeating this millions of times, it gradually learns a good strategy from raw
experience. This is the \textbf{model-free} setting.

\begin{intuition}
``I have no map. Let me just try things. If a move led somewhere good, I'll
remember to do it more often; if it led somewhere bad, I'll do it less.'' That
trial-and-error learning is Q-Learning --- much closer to how an animal (or a
human) actually learns a new environment.
\end{intuition}

\subsection{What is a Q-value?}
Instead of a value per cell, Q-Learning learns a value per
\textbf{(cell, action)} pair, written $Q(s,a)$:
\[
Q(s,a) = \text{``how good is it to take action } a \text{ from cell } s
\text{?''}
\]
If the robot knows $Q(s,a)$ for all actions, it just picks the action with the
largest $Q$. The table of all these numbers is the \textbf{Q-table}.

\subsection{The Q-Learning update (learning from one step)}
After the robot takes action $a$ in cell $s$, lands in cell $s'$, and receives
reward $r$, it nudges its estimate $Q(s,a)$ toward what it just experienced:

\begin{equation}
Q(s,a) \;\leftarrow\; Q(s,a) \;+\; \alpha
\Big[\;\underbrace{r + \gamma \max_{a'} Q(s',a')}_{\text{what I now think it's worth}}
\;-\; \underbrace{Q(s,a)}_{\text{what I thought before}}\;\Big].
\label{eq:qupdate}
\end{equation}

In plain words:
\begin{itemize}[nosep]
  \item $r + \gamma \max_{a'} Q(s',a')$ is the \textbf{TD target}: the reward I
        just got, plus the best I think I can do from the new cell $s'$.
  \item The bracket is the \textbf{error} between my new estimate and my old one.
  \item $\alpha$ (the \textbf{learning rate}, $=0.15$) controls how big a
        correction I make. Small $\alpha$ = learn slowly but steadily.
\end{itemize}

\begin{intuition}
``I expected this move to be worth $X$. It actually turned out to be worth $Y$.
Let me shift my belief a small fraction ($\alpha$) of the way from $X$ toward
$Y$.'' Do this after every single step, for thousands of practice runs, and the
Q-table slowly fills in with accurate values --- \emph{without ever knowing the
map}. It is called \emph{off-policy} because the update always uses
$\max_{a'}$ (the best possible next move) even while the robot is still exploring
randomly.
\end{intuition}

\subsection{Explore vs exploit --- the $\varepsilon$-greedy rule}
A robot that always does what it currently thinks is best will never discover
better routes. So we mix in randomness:

\begin{equation}
\text{action} =
\begin{cases}
\text{a random move} & \text{with probability } \varepsilon \quad(\textbf{explore}),\\
\arg\max_a Q(s,a) & \text{with probability } 1-\varepsilon \quad(\textbf{exploit}).
\end{cases}
\end{equation}

We start with $\varepsilon = 1.0$ (100\% random --- pure exploration, since the
robot knows nothing) and \textbf{decay} it each episode
($\varepsilon \leftarrow 0.999\,\varepsilon$) down to a floor of $0.05$. So early
on the robot explores wildly, and later it mostly follows what it has learned.

\begin{whybox}
Decaying $\varepsilon$ balances the classic \emph{exploration vs exploitation}
dilemma. Early exploration ensures the robot \emph{sees} many routes (including
the good ones); late exploitation ensures it \emph{commits} to the best route it
found. Keeping a small floor ($0.05$) means it never fully stops checking for
something better. This is exactly the ``epsilon decay'' curve in our figures.
\end{whybox}

\subsection{The bitmask trick --- visiting ALL the flowers}
Here is a subtle but crucial problem. The robot must visit \emph{several}
flowers. But an MDP state is memoryless --- from a bare cell, the robot cannot
tell whether it has \emph{already} pollinated flower~3 or not. Without memory it
would loop back to flowers it already visited.

Our fix: \textbf{augment the state} with a small binary memory called a
\emph{bitmask}. With $N=6$ target flowers, we keep a 6-bit number where bit $i$
is $1$ if flower $i$ has been pollinated. The full state becomes:
\[
\text{state} = (\text{row},\; \text{column},\; \text{visited-bitmask}).
\]
So ``cell (4,7) with flowers 0 and 2 already done'' is a \emph{different} state
from ``cell (4,7) with nothing done,'' and the robot can learn different
behaviour for each. The Q-table therefore has shape
$20 \times 20 \times 2^{6} \times 4$.

\begin{intuition}
The bitmask is the robot's tiny checklist: ``flowers done so far.'' Adding it to
the state gives the memoryless robot just enough memory to plan a complete
multi-flower mission and know when it is finally allowed to head to the charger
(only when the checklist is full).
\end{intuition}

\begin{codebox}
In \texttt{QLearning.py}: \texttt{Q = np.zeros((n, n, 1<<N, len(ACTIONS)))} is
the augmented Q-table. \texttt{env\_step(...)} sets a flower's bit with
\texttt{mask | (1 << flower\_bit[nxt])} and only rewards reaching the charger
when \texttt{mask == full\_mask} (all flowers done). \texttt{sample\_move(...)}
is the \emph{only} way the agent touches the world --- it returns just the next
cell and whether it collided, never the transition probabilities.
\end{codebox}

\subsection{Training loop and hyperparameters}
The robot practises for \textbf{6000 episodes}. Each episode: start at the depot
with an empty checklist, take up to \textbf{300 steps} using $\varepsilon$-greedy,
apply the update~\eqref{eq:qupdate} after every step, and stop when the mission
is complete or the step cap is hit. After training we run \textbf{200 greedy test
episodes} (no exploration) to measure the final success rate.

\begin{codebox}
Hyperparameters in \texttt{QLearning.py}: \texttt{ALPHA=0.15} (learning rate),
\texttt{GAMMA=0.95} (discount), \texttt{EPS\_START=1.0}, \texttt{EPS\_MIN=0.05},
\texttt{EPS\_DECAY=0.9990}, \texttt{N\_EPISODES=6000}, \texttt{MAX\_STEPS=300},
\texttt{N\_TEST=200}. A fixed seed makes the whole training reproducible.
\end{codebox}

%=====================================================================
\section{Value Iteration vs Q-Learning --- Side by Side}
%=====================================================================

\begin{table}[h]
\centering
\renewcommand{\arraystretch}{1.3}
\begin{tabular}{@{}p{3.4cm} p{5.0cm} p{5.0cm}@{}}
\toprule
& \textbf{Value Iteration} & \textbf{Q-Learning} \\
\midrule
Knowledge of the world & Knows the full map \& slip model (\textbf{model-based}) & Knows nothing; learns by doing (\textbf{model-free}) \\
How it learns & \emph{Computes} values with the Bellman update & \emph{Learns} $Q$ from trial-and-error experience \\
Needs $T(s,a,s')$? & Yes --- uses it directly & No --- only samples the environment \\
What it stores & Value per cell $U(s)$ & Value per (cell, action) $Q(s,a)$ \\
Exploration? & Not needed (full knowledge) & Essential ($\varepsilon$-greedy) \\
Multi-flower memory & Re-plan per flower (\texttt{run\_mission}) & Bitmask in the state \\
Speed to a solution & Fast ($\sim$36 iterations) & Slower ($\sim$1885 episodes to converge) \\
Result (same slip test) & 100\% success, $\sim$84 steps, reward $58.8$ & 100\% success, $\sim$87 steps, reward $58.4$ \\
Best when\dots & The environment is known & The environment is unknown/changing \\
\bottomrule
\end{tabular}
\caption{Two ways to solve the identical navigation MDP.}
\end{table}

\subsection{Why we implement both}
They answer the same navigation question under \emph{opposite assumptions},
which is the whole pedagogical point:
\begin{itemize}[nosep]
  \item \textbf{Value Iteration} is the ``gold standard'' when you have a perfect
        model. It is fast and gives the provably optimal policy --- a great
        \emph{benchmark}.
  \item \textbf{Q-Learning} needs no model, so it works in realistic settings
        where a robot is dropped into an \emph{unknown} greenhouse. It is slower
        and noisier, but far more flexible.
\end{itemize}
Because they share the \emph{identical} grid, obstacles, flowers and rewards
(Q-Learning imports the \texttt{GridWorld} from \texttt{ValueIteration.py}), we
can fairly compare ``planning with a model'' against ``learning without one.''
To keep the comparison honest, \textbf{both} policies are scored the same way ---
rolled out over 200 episodes \emph{with the $0.8/0.1/0.1$ actuator noise
switched on} (function \texttt{evaluate\_vi\_policy} does this for Value
Iteration, exactly mirroring Q-Learning's greedy test). Under that identical
test, Q-Learning rediscovers almost the same quality of route (reward $58.4$ vs
$58.8$, $\sim$87 vs $\sim$84 steps) that Value Iteration computes instantly ---
a near-optimal result reached with \emph{no map at all}.

\begin{tcolorbox}[colback=green!5!white,colframe=green!45!black,title=Take-away]
Both are \textbf{reinforcement learning} on the same MDP. Value Iteration
\emph{plans} with a known model (fast, exact, needs the map). Q-Learning
\emph{learns} from experience with no model (flexible, slower, needs
exploration). Having both demonstrates the fundamental \emph{model-based vs
model-free} trade-off at the centre of reinforcement learning.
\end{tcolorbox}

%=====================================================================
\section{Where This Fits in the Whole Project}
%=====================================================================

\begin{center}
\begin{tabular}{@{}ll@{}}
\toprule
\textbf{Phase} & \textbf{Job} \\
\midrule
Phase 1 & PSO / GA decide \emph{which robot} gets \emph{which flowers}. \\
Phase 2 (this document) & Value Iteration \& Q-Learning navigate the robot \\
                        & cell-by-cell to actually reach those flowers. \\
\bottomrule
\end{tabular}
\end{center}

\noindent
Phase~1 answered ``\textbf{who does what}.'' Phase~2 answers ``\textbf{how do I
get there},'' taking the flower assignments as fixed input. The two phases never
mix algorithms: population-based optimization handles allocation, reinforcement
learning handles navigation --- each staying strictly in its own lane.

\end{document}
